% ============================================================
% 片段：s03b  对应原文 PDF 第 18--22 页
% 内容：§3.5 B 模型（3.5.1 B 扭转、3.5.2 霍奇结构的变化、3.5.3 特殊/平坦坐标）
% 译者：t-s03b   审校状态：初稿
% ============================================================

% ==== 原文 PDF 第 18 页 ====

\subsection{B 模型}
\label{sec:3.5}

\begin{keytakeaways}
  \begin{itemize}[leftmargin=1.4em,itemsep=2pt,topsep=2pt]
    \item \textbf{物理图像}：B 模型是镜像对称中“复结构一侧”的拓扑弦。B 扭转后，世界面映射的鞍点条件迫使映射为常值映射，因此 B 模型不含瞬子修正，其关联函数是靶空间几何量（度规、曲率、全纯 $(3,0)$ 形式）的经典积分。
    \item \textbf{数学工具}：B 模型的上同调由靶空间 $Z^*$ 的 Dolbeault 上同调实现：算符与取值于 $\wedge^pTZ^*$ 的 $(0,q)$ 形式一一对应，BRST 算符 $Q_B$ 对应 $\bar\partial$。
    \item \textbf{关键结论}：荷为 $(1,1)$ 的形变算符经全纯 $(3,0)$ 形式 $\Omega$ 缩并后成为 $(2,1)$ 形式，故形变子环由 $H^{2,1}(Z^*)$ 承载，它参数化 $Z^*$ 的复结构形变。B 模型侧的计算因此归结为霍奇结构的变化与周期积分。
  \end{itemize}
\end{keytakeaways}

% [Tutor Note: 物理图像] B 模型为何“经典”而无瞬子修正
\begin{yizhu}[物理图像]
B 模型与 A 模型的本质差别在于世界面映射的鞍点方程。令 (3.27) 中的 $\mathcal{D}=0$ 后，$\partial\phi^{\bar k}=\bar\partial\phi^{\bar k}=0$，即映射 $\phi:\Sigma\to Z^*$ 必须是常值映射，世界面塌缩为靶空间中的一点。瞬子修正来自绕数非零的全纯映射（即 A 模型所允许的非常值解），故 B 模型完全没有此类贡献，其关联函数退化为经典积分。这也解释了为什么 B 模型只探测复结构模（经由 $\Omega$ 的周期），而 A 模型探测凯勒模（经由世界面瞬子的计数）。
\end{yizhu}

% [Tutor Note: 易错点] 手征环的“荷 (1,1)”与 H^{2,1}(Z^*) 的对应
\begin{yizhu}[易错点]
“荷为 $(1,1)$ 的形变算符对应 $H^{2,1}(Z^*)$ 的元素”一步靠的是用 $\Omega$ 作缩并：由 (3.24)--(3.25)，算符由取值于 $\wedge^pTZ^*$ 的 $(0,q)$ 形式刻画，其中 $p$ 来自 $\theta_{i_1}\cdots\theta_{i_p}$（全纯切向指标），$q$ 来自 $\eta^{\bar\imath_1}\cdots\eta^{\bar\imath_q}$（反全纯微分指标）。与全纯 $(3,0)$ 形式 $\Omega$ 缩并后变为 $(3-p,q)$ 形式，于是 $(p,q)=(1,1)$ 给出 $(3-1,1)=(2,1)$。核对 (3.21) 的丛分配可确认这两个指标的来源不同，切勿把 $p,q$ 与时空荷混为一谈。
\end{yizhu}

% [Tutor Note: 易错点] 靶空间记号 Z 与 Z^* 的使用
\begin{yizhu}[易错点]
原文在 (3.21) 中把靶空间写作 $Z$，而自 (3.24) 之后一律写作 $Z^*$；按 §3.5.2 开头“B 侧的模空间对应靶空间 $Z^*$ 的复结构模空间”的说法，本节 B 模型的靶空间应是 $Z^*$（$Z$ 是 A 模型一侧的 CY 流形）。引用指标时须注意 $Z$ 与 $Z^*$ 指同一族几何对象在两侧的成员，不是两个无关空间。
\end{yizhu}

\subsubsection{B 扭转}
\label{sec:3.5.1}

\begin{keytakeaways}
  \begin{itemize}[leftmargin=1.4em,itemsep=2pt,topsep=2pt]
    \item \textbf{物理图像}：B 扭转把四个费米子重新分组为一个 1-形式 $\rho^i$（由 $\psi^i_z$ 与 $\chi^i_{\bar z}$ 组成）和两个标量 $\eta^{\bar\imath}$、$\theta_j$，从而使 BRST 算符成为纯反全纯微分。
    \item \textbf{数学工具}：$\eta^{\bar\imath}\leftrightarrow \mathrm{d}\bar u^{\bar\imath}$、$\theta_i\leftrightarrow \partial/\partial u^i$、$Q_B\leftrightarrow \bar\partial$，于是拓扑算符构成取值于 $\wedge^pTZ^*$ 的 $(0,q)$ 形式的 de Rham--Dolbeault 复形。
    \item \textbf{关键结论}：作用量可写成 $S=i\int_\Sigma\{Q_B,\mathcal{D}\}+U$ 的 $Q_B$-恰当形式加经典项 $U$；$Q_B$-闭模去 $Q_B$-恰当类即给出 Dolbeault 上同调，这正是 B 模型的手征环。
  \end{itemize}
\end{keytakeaways}

回到 B 扭转（3.2），费米子作为其截面的丛的改动为
\begin{equation}\tag{3.21}
\begin{aligned}
\psi^i_z &\in \Gamma(K\otimes\phi^*T^{1,0}Z)\,, &\qquad
\psi^{\bar\imath} &\in \Gamma(\phi^*T^{0,1}Z)\,,\\
\chi^i_{\bar z} &\in \Gamma(\bar K\otimes\phi^*T^{1,0}Z)\,, &\qquad
\chi^{\bar\imath} &\in \Gamma(\phi^*T^{0,1}Z)\,.
\end{aligned}
\end{equation}

% [Tutor Note: 易错点] (3.21) 中 K、\bar K 与 z、\bar z 的配对
\begin{yizhu}[易错点]
(3.21) 的关键是“丛与下标要配对”：$\psi^i_z$ 的下标 $z$ 配典范丛 $K$，$\chi^i_{\bar z}$ 的下标 $\bar z$ 配反典范丛 $\bar K$，二者都取全纯切丛 $\phi^*T^{1,0}Z$ 的值；而 $\psi^{\bar\imath}$、$\chi^{\bar\imath}$ 不带世界面下标（即 $(0,0)$ 形式），取反全纯切丛 $\phi^*T^{0,1}Z$ 的值。与 A 扭转相比，(3.21) 的“错位”正是 B 扭转把 $(0,1)$ 世界面指标从 $\bar K$ 上去掉的结果，它使 $\eta,\theta$ 成为世界面标量，从而是 $Q_B=\bar\partial$ 而非 $Q_A=\partial$ 的候选。
\end{yizhu}

世界面标量按如下方式分组：
\begin{equation}\tag{3.22}
\eta^{\bar\imath}=\psi^{\bar\imath}+\chi^{\bar\imath}\,,\qquad
\theta_j=g_{\bar\imath j}\bigl(\psi^{\bar\imath}-\chi^{\bar\imath}\bigr)\,,
\end{equation}
而 $\psi^i_z,\chi^i_z$ 被视为一个 1-形式 $\rho^i$ 的 $(1,0)$ 与 $(0,1)$ 分量。
% [原文疑误: 本句中的 \chi^i_z 应为 \chi^i_{\bar z}，与上一式 (3.21) 中 \chi^i_{\bar z}\in\Gamma(\bar K\otimes\phi^*T^{1,0}Z) 以及“\rho^i 的 (1,0) 与 (0,1) 分量”的说法一致；原文排为 \chi^i_z，照抄于此。]
% [Tutor Note: 原文疑误] (3.22) 后一句的 \chi^i_z 与 (3.21) 的 \chi^i_{\bar z} 不一致
\begin{yizhu}[原文疑误]
紧随 (3.22) 的那句“$\psi^i_z,\chi^i_z$ 是一个 1-形式 $\rho^i$ 的 $(1,0)$ 与 $(0,1)$ 分量”中，第二个下标原文排作 $z$，而 (3.21) 中相应对象是 $\chi^i_{\bar z}\in\Gamma(\bar K\otimes\phi^*T^{1,0}Z)$。既然 $\rho^i$ 的 $(0,1)$ 分量应带 $\mathrm{d}\bar z$，此处应为 $\chi^i_{\bar z}$；译文照抄原文的 $z$，读时按 $\bar z$ 理解即可。
\end{yizhu}

为分析
\[
Q_B=Q^++\bar Q^+,
\]
的上同调，在 (2.32-2.37) 中取 $\varepsilon^+=\bar\varepsilon^+=\varepsilon$ 以及 $\bar\varepsilon^+=\varepsilon^+=0$，使得超对称变分 $\delta$ 由
\[
\delta=\varepsilon\,Q_B
\]
给出。
% [原文疑误: 原文第二组条件排作 \bar\varepsilon^+=\varepsilon^+=0，与紧接其前的 \varepsilon^+=\bar\varepsilon^+=\varepsilon 直接矛盾（同一分量不能既等于 \varepsilon 又等于 0）。按上下文 (2.32--2.37) 中的 \varepsilon^{\pm} 两个分量，第二组应为 \varepsilon^-=\bar\varepsilon^-=0。照抄原文的上标 +。]
% [Tutor Note: 原文疑误] 第二组超对称参数条件的上标应为负号
\begin{yizhu}[原文疑误]
“取 $\varepsilon^+=\bar\varepsilon^+=\varepsilon$ 以及 $\bar\varepsilon^+=\varepsilon^+=0$”两句无法同时成立：若 $\varepsilon^+=\bar\varepsilon^+=\varepsilon$，则同一分量不能又等于 $0$。参照 (2.32--2.37) 中 $\varepsilon^{\pm}$ 表示两个独立超对称参数，第二组条件应为 $\varepsilon^-=\bar\varepsilon^-=0$（即只保留使 $Q_B$ 成为 BRST 算符的那一半超对称）。译文照抄原文的上标，特此标出。
\end{yizhu}

超对称变分提示如下对应 \cite{ref6}：
\begin{equation}\tag{3.23}
\eta^{\bar\imath}\leftrightarrow \mathrm{d}\bar u^{\bar\imath}\,,\qquad
\theta_i\leftrightarrow\frac{\partial}{\partial u^i}\,,\qquad
Q_B\leftrightarrow\bar\partial\,,
\end{equation}

由二维标量构成的拓扑场论算符具有形式
\begin{equation}\tag{3.24}
\omega(\phi)^{i_1,\ldots,i_p}_{\bar\imath_1,\ldots,\bar\imath_q}\,
\eta^{\bar\imath_1}\ldots\eta^{\bar\imath_q}\,
\theta_{i_1}\ldots\theta_{i_p}\,,
\end{equation}
它们于是可以等同于
\begin{equation}\tag{3.25}
\mathrm{d}\bar u^{\bar\imath_1}\ldots \mathrm{d}\bar u^{\bar\imath_q}\;
\omega(u)^{i_1,\ldots,i_p}_{\bar\imath_1,\ldots,\bar\imath_q}\;
\frac{\partial}{\partial u^{i_1}}\ldots\frac{\partial}{\partial u^{i_p}}\,.
\end{equation}

拓扑理论的算符与 $Z^*$ 上取值于 $\wedge^pTZ^*$ 的 $(0,q)$ 形式一一对应。进一步，作用量现在可以写成
\begin{equation}\tag{3.26}
S=i\int_\Sigma\{Q_B,\mathcal{D}\}+U\,,
\end{equation}
其中
\begin{equation}\tag{3.27}
\mathcal{D}=g_{j\bar k}\Bigl(\rho^j_z\,\bar\partial\phi^{\bar k}+\rho^j_{\bar z}\,\partial\phi^{\bar k}\Bigr)\,,
\end{equation}
\begin{equation}\tag{3.28}
U=\int_\Sigma\Bigl(-\theta_j D\rho^j-\frac{i}{2}R_{j\bar\jmath k\bar k}\,
\rho^j\wedge\rho^k\,\eta^{\bar j}\theta_l\,g^{l\bar k}\Bigr)\,.
\end{equation}
% [原文疑误: (3.28) 括号内第一项排作普通斜体 D\rho^j，而 (3.26)、(3.27) 以及第 19 页开头的 "Setting \mathcal{D}=0" 均用花体 \mathcal{D}；按上下文二者是同一算符，原文此处字体不一致，照抄为 D。]
% [Tutor Note: 易错点] (3.28) 的 D 与 \mathcal{D}、R 的指标顺序与带点/无点 j
\begin{yizhu}[易错点]
三处易错：(i) (3.28) 第一项原文用普通 $D$，而 (3.26)、(3.27) 与第 19 页“令 $\mathcal{D}=0$”用花体 $\mathcal{D}$，应读作同一辅助场；(ii) 曲率张量的指标顺序是 $R_{j\bar\jmath k\bar k}$，与因子 $\rho^j\wedge\rho^k\,\eta^{\bar j}\theta_l\,g^{l\bar k}$ 按 $j,\bar\jmath,k,\bar k$ 一一配对，升降指标全部用度规 $g$、$g^{-1}$，核对此式时不要改动指标次序；(iii) 原文中 $R$ 的下标 $\bar\jmath$ 排为无点 $j$，而同一式的 $\eta^{\bar j}$ 排为带点 $j$，字体不一致，照抄原文。
\end{yizhu}

% ==== 原文 PDF 第 19 页 ====

令 $\mathcal{D}=0$，映射 $\phi$ 变得平凡，因为 $\partial\phi^{\bar k}=\bar\partial\phi^{\bar k}=0$；这些映射把二维世界面映为 $Z^*$ 中的一点。因此 B 模型不接收任何瞬子修正。关联函数对应于经典积分。此外，CY 三维流形 $Z^*$ 的全纯 $(3,0)$ 形式 $\Omega$（在相差一个常数意义下唯一）可用来把取值于 $\wedge^pTZ^*$ 的 $(0,q)$ 形式变换为 $(3-p,q)$ 形式。荷为 $(1,1)$ 的形变算符因此被等同于 $H^{2,1}(Z^*)$ 中的元素，而后者参数化复结构形变。

\subsubsection{霍奇结构的变化}
\label{sec:3.5.2}

\begin{keytakeaways}
  \begin{itemize}[leftmargin=1.4em,itemsep=2pt,topsep=2pt]
    \item \textbf{物理图像}：B 侧的模空间就是靶空间 $Z^*$ 的复结构模空间；形变子环由中间维上同调 $H^3(Z^*,\mathbb{C})$ 承载，而 $H^3$ 的霍奇分解把复结构形变的信息组织成“荷”的层级。
    \item \textbf{数学工具}：$H^{p,q}(Z^*)$ 分解不是全纯的，但霍奇滤过 $F^p(Z^*)=\bigoplus_{a\ge p}H^{a,3-a}(Z^*)$ 在全纯变化下稳定，且 $H^{p,q}=F^p\cap\overline{F^q}$ 可还原原分解。
    \item \textbf{关键结论}：Gauss--Manin 联络 $\nabla$ 满足 Griffiths 横截性 $\nabla F^p\subset F^{p-1}$，据此 $\Omega(x)$ 的多重导数逐级填满整个滤过；周期被四阶的 Picard--Fuchs 方程零化，该方程即霍奇结构变化的微分方程表达。三点函数 $C_{ijk}$ 由 $\Omega\wedge\nabla_i\nabla_j\nabla_k\Omega$ 的积分给出。
  \end{itemize}
\end{keytakeaways}

B 侧的模空间对应于靶空间 $Z^*$ 的复结构模空间。对 CY 三维流形，形变子环所对应的丛是中间维上同调 $H^3(Z^*,\mathbb{C})$。一旦选定某个复结构，即选定复结构模空间中的某个特定点，该空间便有一个自然的分解。分解为
\begin{equation}\tag{3.29}
H^3(Z^*,\mathbb{C})\simeq\bigoplus_{p+q=3}H^{p,q}(Z^*)\,.
\end{equation}
该分解确定了全纯 $(3,0)$ 形式 $\Omega$（在相差一个常数意义下唯一）。它还给出复共轭的自然概念，即 $\overline{H^{p,q}(Y)}=H^{q,p}(Y)$。术语“霍奇结构”（Hodge structure）指 $H^3(Z^*,\mathbb{C})$ 连同分解 (3.29) 以及一个由 $H^3(Z^*,\mathbb{Z})$ 给出的格，后者与 $\mathbb{C}$ 张量后生成 $H^3(Z^*,\mathbb{C})$（详见文献 \cite{ref48,ref49} 及其中的参考文献）。然而当复结构模被改变时，分解 (3.29) 并不全纯地变化。这一分解的非全纯性正是后文将要讨论的全纯反常方程的根源。
% [原文疑误: 此处及 (3.35) 中共轭公式的流形记作 Y，但本节靶空间一律记作 Z^*，全文未定义 Y；按上下文 Y 即 Z^*。照抄原文的 Y。]

不过，存在一种\emph{不同}的丛分解，它在复结构模空间上全纯地变化。该分解由霍奇滤过（Hodge filtration）$F^{\bullet}(Z^*)=\{F^p(Z^*)\}_{p=0}^{3}$ 给出，其中括号内的空间由下式定义
\begin{equation}\tag{3.30}
H^3=F^0\supset F^1\supset F^2\supset F^3\supset F^4=0\,,\qquad
F^p(Z^*)=\bigoplus_{a\ge p}H^{a,3-a}(Z^*)\subset H^3\,.
\end{equation}
为恢复分解 (3.29)，可将霍奇滤过与反全纯滤过相交
\begin{equation}\tag{3.31}
H^{p,q}(Z^*)=F^p(Z^*)\cap\overline{F^q(Z^*)}\,.
\end{equation}
% [Tutor Note: 易错点] F^p 的维数与手征环荷的对应
\begin{yizhu}[易错点]
霍奇滤过与 $(p,q)$ 分解的对应是：$F^p=\bigoplus_{a\ge p}H^{a,3-a}$，故对 CY 三维流形 $\dim F^3=1$（由 $\Omega$ 张成）、$\dim F^2=1+h^{2,1}$、$\dim F^1=1+2h^{2,1}$、$\dim F^0=2+2h^{2,1}=\dim H^3$。把它与手征环的荷分层对照时，$F^p$ 的指标 $p$ 是“全纯形式次数”，插入一个单位荷的边缘算符恰好把 $p$ 降 $1$（见 Griffiths 横截性），这与 (3.44) 中联络矩阵的上三角结构是同一件事的两种说法。另外 (3.31) 的共轭是对 $F^q$ 整体取，不要漏掉上划线。
\end{yizhu}

% [Tutor Note: 含糊点] “分解不全纯地变化”的实质
\begin{yizhu}[含糊点]
原文只说 (3.29) 的分解“当复结构模被改变时并不全纯地变化”，其被省略的实质是：$H^{p,q}$ 子空间依赖于模空间上的点，作为模空间上的丛并不分解为全纯子丛；能全纯变化的只有滤过 $F^p$（这正是 (3.30) 的意义），而 $F^p$ 与 $\overline{F^q}$ 的交才还原出 $H^{p,q}$。也正因如此，周期向量 $w(x)$ 可以全纯地定义（(3.33)）；而把 $H^{p,q}$ 直接当作模空间上的全纯对象来用，正是后文全纯反常方程所处理的困难的来源。
\end{yizhu}

代替固定的靶空间，可以通过改变 $Z^*$ 的复结构来考虑一个形变族。此时滤过上配备一个平坦联络 $\nabla$，称为 Gauss--Manin 联络，其性质 $\nabla F^p\subset F^{p-1}$ 称为 Griffiths 横截性。
这一性质使得 $\Omega(x)$\footnote{$x$ 是复结构模空间 $\mathcal{M}$ 上的局部坐标，用来标明形式类型对模空间中所处坐标卡的依赖。} 的导数（$F^3$ 中的元素）可以被等同于较低滤过空间中的元素。整个滤过可以通过取全纯 $(3,0)$ 形式的多重导数来张成。

% ==== 原文 PDF 第 20 页 ====

四阶导数于是又可以由基的元素表示，这体现在 $\Omega(x)$ 的周期被一个四阶微分方程组所零化，该方程组称为 Picard--Fuchs（PF）方程。PF 方程刻画霍奇结构的变化，而后者在 B 模型一侧给出了 $\mathcal{N}=(2,2)$ 超共形场论及其手征环 \cite{ref50} 的形变的几何实现；综述另见文献 \cite{ref51}。示意地，这一变化具有如下形式：
\begin{equation}\tag{3.32}
F^3\xrightarrow{\nabla_B}F^2\xrightarrow{\nabla_B}F^1\xrightarrow{\nabla_B}F^0\,.
\end{equation}
空间 $(F^3,\,F^2/F^3,\,F^1/F^2,\,F^0/F^1)$ 的维数依次为 $(1,h^{2,1},h^{2,1},1)$。这些空间中的元素可通过对 $\Omega(x)$ 关于模求导得到。滤过的一个截面可由如下向量 $w(x)$ 给出，它有 $2h^{2,1}+2$ 个分量：
\begin{equation}\tag{3.33}
w(x)=\bigl(\Omega(x)\quad \theta_i\Omega(x)\quad \theta_*\theta_i\Omega(x)\quad \theta_*^3\Omega(x)\bigr)^t\,.
\end{equation}
其中 $\theta_i=x^i\partial_{x^i}$，$i=1\ldots,h^{2,1}(Z^*)$，$x^i$ 为某些局部坐标，而 $x^*$ 是某个固定的坐标选取，使得 $h^{2,1}$ 个元素 $\theta_i\theta_*\Omega(x)$ 张成 $F^1/F^2$。利用 $w(x)$ 可定义周期矩阵
\begin{equation}\tag{3.34}
\Pi(x)^{\alpha}_{\beta}=\int_{\gamma^\alpha}w(x)\,,\qquad
\gamma^\alpha\in H_3(Z^*)\,,\qquad
\alpha,\beta=0,\ldots,2h^{2,1}+1\,,
\end{equation}
其第一行对应于 $\Omega(x)$ 的周期。周期被 PF 算符零化。

% [Tutor Note: 含糊点] Picard--Fuchs 方程的来源
\begin{yizhu}[含糊点]
原文说“四阶导数又可以由基的元素表示”，其省略的推理是：由 Griffiths 横截性，对模空间方向的逐次求导（以 $\partial$ 记之）有 $\Omega\in F^3$、$\partial\Omega\in F^2$、$\partial^2\Omega\in F^1$、$\partial^3\Omega\in F^0$，而 $\dim H^3=2h^{2,1}+2$ 有限，故足够高阶的导数必是前面诸元素的线性组合。对一参数形变族（如 §3.6 的五次超曲面，$h^{2,1}=1$）此时恰好是第 $4$ 阶，于是每个周期 $\int_\gamma\Omega$ 满足一个四阶常微分方程，即 PF 方程；一般 $h^{2,1}>1$ 时得到的是 PF 方程组（Pfaffian 系统），原文的“a system of differential equations”即指这种一般情形，而阶数随具体族而定。
\end{yizhu}

此外，Griffiths 横截性连同由在 CY 上对两个 3-形式的楔积作积分所给出的辛配对，使人可以定义三点函数（three-point functions）\footnote{此处的符号是一个约定选择。}
\begin{equation}\tag{3.35}
C_{ijk}:=-\int_Y\Omega\wedge\nabla_i\nabla_j\nabla_k\Omega\,.
\end{equation}
% [Tutor Note: 原文疑误] (3.35) 的积分域写作 Y
\begin{yizhu}[原文疑误]
(3.35) 的积分域原文写作 $Y$（(3.31) 中的共轭公式也写作 $Y$），而本节其余各处靶空间均记作 $Z^*$，且全文未引入 $Y$。按上下文 $Y$ 即 $Z^*$。译文照抄 $Y$ 以便与原文逐字核对；同时注意它与 (3.42) 中以 $Z^*$ 为积分域的汤川耦合定义只差归一化以及 $\Omega$ 与 $v_0$ 的取法。
\end{yizhu}

% ==== 原文 PDF 第 21 页 ====

\subsubsection{特殊/平坦坐标}
\label{sec:3.5.3}

\begin{keytakeaways}
  \begin{itemize}[leftmargin=1.4em,itemsep=2pt,topsep=2pt]
    \item \textbf{物理图像}：特殊坐标 $t^a=X^a/X^0$ 把复结构模空间上的几何量与底层 $\mathcal{N}=2$ 理论的物理形变参数对应起来，因而正是镜像对称中的镜像映射。
    \item \textbf{数学工具}：在辛基 $A^I,B_J$ 与对偶基 $\alpha_I,\beta^J$ 下展开 $\Omega=X^I\alpha_I+\mathcal{F}_J\beta^J$；$X^I$ 给出射影坐标，$\mathcal{F}_J$ 是权为 2 的齐次函数（预势）$F_0(t)$ 的导数。
    \item \textbf{关键结论}：在特殊坐标下，Gauss--Manin 联络取 (3.45) 的形式 $(\partial_a-C_a)v=0$，其中 $C_a$ 由上三角矩阵给出；上三角性正是 Griffiths 横截性（等价于插入单位荷边缘算符使手征环元素荷加 $1$）的矩阵表述。
    \item \textbf{计算要点}：矩阵的非零元为 $\delta^c_a$、$C_{abc}$、$\delta^b_a$，汤川耦合 $C_{abc}=\partial_a\partial_b\partial_cF_0(t)$ 因此就是预势的三阶导数。
  \end{itemize}
\end{keytakeaways}

下面讨论复结构模空间上的一组特殊坐标，它们使得与底层理论的物理形变建立对应成为可能。这组坐标定义于特殊几何（special geometry）之内，而特殊几何正是在研究 $\mathcal{N}=2$ 理论的模空间过程中发展起来的。\footnote{参见文献 \cite{ref52,ref53,ref54,ref5,ref55,ref12,ref56,ref50,ref57,ref58,ref51,ref59}。}
选取 3-闭链的辛基 $A^I,B_J\in H_3(Z^*)$ 以及 $H_3(Z^*)$ 的一组对偶基 $\alpha_I,\beta^J$，使得
\begin{equation}\tag{3.36}
\begin{aligned}
&A^I\cap B_J=\delta^I_{\ J}=-B_J\cap A^I\,,\qquad A^I\cap A^J=B_I\cap B_J=0\,,\\
&\int_{A^I}\alpha_J=\delta^I_J\,,\qquad \int_{B_J}\beta^I=\delta^I_J\,,\qquad I,J=0,\ldots,h^{2,1}(Z^*)\,,
\end{aligned}
\end{equation}
% [Tutor Note: 易错点] 辛基、对偶基与脚注 14 的符号约定
\begin{yizhu}[易错点]
(3.36) 两组条件要分开读：第一行是 3-闭链的交截数（斜对称，故 $A^I\cap B_J=-B_J\cap A^I$），第二行是 3-形式在闭链上的周期（$\alpha_J$ 在 $A^I$ 上归一、$\beta^I$ 在 $B_J$ 上归一）。指标位置很要紧：$A^I$ 带上标、$B_J$ 带下标，因此 $\delta^I_{\ J}$ 的第二个指标要降低。脚注 14 说明了与之配套的符号约定 $X^I=\int_{A^I}\Omega$、$F_I=\int_{B_I}\Omega$，恰与 (3.37) 的展开 $\Omega=X^I\alpha_I+\mathcal{F}_J\beta^J$ 自洽；换用相反约定会让下面 (3.38)--(3.41) 的若干符号整体反号。
\end{yizhu}

全纯 $(3,0)$ 形式 $\Omega(x)$ 可在基 $\alpha_I,\beta^J$ 中展开：\footnote{相对符号的选取是一个约定，文献中各不相同；此处选取使 $X^I=\int_{A^I}\Omega$、$F_I=\int_{B_I}\Omega$。}
\begin{equation}\tag{3.37}
\Omega(x)=X^I(x)\alpha_I+\mathcal{F}_J(x)\beta^J\,.
\end{equation}
周期 $X^I(x)$ 可以等同于 $\mathcal{M}$ 上的射影坐标，而 $\mathcal{F}_J$ 等同于一个权为 2 的齐次函数 $\mathcal{F}(X^I)$ 的导数，即 $\mathcal{F}_J=\partial\mathcal{F}(X^I)/\partial X^J$。在 $X^0(x)\neq 0$ 的坐标卡中，可以定义一组特殊坐标
\[
t^a=\frac{X^a}{X^0}\,,\qquad a=1,\ldots,h^{2,1}(Z^*)\,.
\]
归一化的全纯 $(3,0)$ 形式 $v_0=(X^0)^{-1}\Omega(t)$ 具有展开
\begin{equation}\tag{3.38}
v_0=\alpha_0+t^a\alpha_a+\beta^bF_b(t)+\bigl(2F_0(t)-t^cF_c(t)\bigr)\beta^0\,,
\end{equation}
其中
\[
F_0(t)=(X^0)^{-2}\mathcal{F}\,,\qquad
F_a(t):=\partial_aF_0(t)=\frac{\partial F_0(t)}{\partial t^a}\,.
\]
$F_0(t)$ 是预势（prepotential）。可以定义
\begin{align}
v_a&=\alpha_a+\beta^bF_{ab}(t)+\bigl(F_a(t)-t^bF_{ab}(t)\bigr)\beta^0\,, \tag{3.39}\\
v^a_D&=\beta^a-t^a\beta^0\,, \tag{3.40}\\
v^0&=-\beta^0\,. \tag{3.41}
\end{align}
% [Tutor Note: 物理图像] t^a=X^a/X^0 为什么使联络上三角
\begin{yizhu}[物理图像]
$t^a=X^a/X^0$ 之所以称为“特殊/平坦坐标”，是因为在对这些坐标求导时，Gauss--Manin 联络的矩阵逐级升高一阶：$\partial_a$ 作用在 $v_0\in F^3$ 上给出 $v_a\in F^2$（对应 $\delta^c_a$），作用在 $v_b\in F^2$ 上给出 $C_{abc}v^c_D\in F^1$，作用在 $v^b_D\in F^1$ 上给出 $\delta^b_a v^0\in F^0$，而 $\partial_a v^0=0$。这正是 Griffiths 横截性 $\nabla F^p\subset F^{p-1}$ 的矩阵形式，其物理含义是：插入一个单位荷的边缘算符使手征环元素的荷恰好加 $1$。在一般（非特殊）坐标下这一结构会被非平凡的联络系数搅乱；特殊坐标把它化为 (3.44) 的上三角常矩阵，故这些坐标就是与 A 模型平坦坐标相认同的镜像映射。
\end{yizhu}

特殊坐标下的汤川耦合（Yukawa coupling）由下式给出
\begin{equation}\tag{3.42}
C_{abc}:=\partial_a\partial_b\partial_cF_0(t)=-\int_{Z^*}v_0\wedge\partial_a\partial_b\partial_cv_0\,.
\end{equation}
进一步定义具有 $2h^{2,1}+2$ 个分量的向量
\begin{equation}\tag{3.43}
v=\bigl(v_0\,,\quad v_a\,,\quad v^a_D\,,\quad v^0\bigr)^t\,,
\end{equation}
按构造它给出
\begin{equation}\tag{3.44}
\partial_a
\begin{pmatrix}
v_0\\ v_b\\ v^b_D\\ v^0
\end{pmatrix}
=\underbrace{
\begin{pmatrix}
0 & \delta^c_a & 0 & 0\\
0 & 0 & C_{abc} & 0\\
0 & 0 & 0 & \delta^b_a\\
0 & 0 & 0 & 0
\end{pmatrix}}_{:=C_a}
\begin{pmatrix}
v_0\\ v_c\\ v^c_D\\ v^0
\end{pmatrix},
\end{equation}
% [Tutor Note: 易错点] (3.44) 中矩阵元的位置与分量顺序
\begin{yizhu}[易错点]
读 (3.44) 时必须把 $v$ 的分量顺序 $v=(v_0,\,v_b,\,v^b_D,\,v^0)^t$ 与矩阵的列顺序对齐：第 $1$ 行第 $2$ 列的 $\delta^c_a$ 把 $v_0$ 送到 $v_a$；第 $2$ 行第 $3$ 列的 $C_{abc}$ 把 $v_b$ 送到 $C_{abc}v^c_D$；第 $3$ 行第 $4$ 列的 $\delta^b_a$ 把 $v^b_D$ 送到 $\delta^b_a v^0$（注意 $v^0=-\beta^0$ 已含符号，故 $\partial_a(\beta^b-t^b\beta^0)=-\delta^b_a\beta^0=\delta^b_a v^0$）；第 $4$ 行全为零，即 $\partial_av^0=0$。若把 $v$ 的分量换序（例如把 $v^a_D$ 与 $v_b$ 对调），矩阵形状随之改变，切勿照抄符号而不核对分量顺序。
\end{yizhu}

由此定义了 $(2h^{2,1}+2)\times(2h^{2,1}+2)$ 矩阵 $C_a$，借助它们该方程可写成如下形式：
\begin{equation}\tag{3.45}
(\partial_a-C_a)\,v=0\,.
\end{equation}

% ==== 原文 PDF 第 22 页 ====

$v$ 的各分量对应于前面讨论的各个滤过空间中的元素。式 (3.45) 在特殊坐标下定义了 Gauss--Manin 联络。联络矩阵的上三角结构反映了：在手征环中插入一个单位荷的边缘算符时，其中元素的荷增加所带来的效应。由于底层超共形场论对 A 模型与 B 模型是同构的，这套描述霍奇结构变化的坐标正是描述镜像对称的合适坐标，并因此给出镜像映射。镜像对称的力量来自这样一个事实：B 模型一侧的计算因与复结构变化相联系而得到很好的控制；所涉及的计算是经典的。与 A 模型非平凡量子几何的对应，是通过找出正确的镜像映射（mirror map）集合来实现的，而这些映射就是特殊坐标，也称为平坦坐标（flat coordinate）。它们正是使 Gauss--Manin 联络取式 (3.44) 形式的坐标。有若干技术 \cite{ref50,ref57,ref58,ref51,ref60,ref61,ref62,ref63,ref28} 可用来确定这一特殊坐标集，从而与物理参数、进而与 A 模型一侧建立对应。下文将以五次超曲面为例对此作一概述。

% [Tutor Note: 含糊点] “若干技术”指什么
\begin{yizhu}[含糊点]
原文以“有若干技术 \cite{ref50,ref57,ref58,ref51,ref60,ref61,ref62,ref63,ref28} 可用来确定特殊坐标集”一句带过具体做法，被省略的是：先由 Picard--Fuchs 算符定出周期解（即 (3.34) 的 $\Pi$），再取周期的比值 $t^a=X^a/X^0$ 作为模空间的平坦坐标，并验证在这些坐标下 Gauss--Manin 联络的矩阵恰为上三角的常矩阵 $C_a$。原文随后以五次超曲面为例说明该流程（见 §3.6），此处不宜凭空补出任何具体变换公式。
\end{yizhu}
